\chapter{Jerarquía mínima}

\section{Noción de jerarquía}
\begin{definition}
Se llama jerarquía mínima al orden parcial de niveles de tratamiento dentro del canon, determinado por dependencia estructural, ganancia operativa y capacidad de organización creciente.
\end{definition}

\begin{observation}
La jerarquía mínima no es todavía una escala exhaustiva del sistema. Su función es declarar una relación ordenada entre formulaciones sin colapsarlas en un único plano indiferenciado.
\end{observation}

\section{Relación de precedencia}
\begin{definition}
Se dice que una formulación \(A\) precede jerárquicamente a una formulación \(B\) cuando \(B\) depende de \(A\), la conserva como herencia canónica y añade una capacidad operativa o estructural que \(A\) no poseía en el mismo nivel.
\end{definition}

\begin{observation}
La precedencia jerárquica no implica superioridad absoluta. Solo fija una dependencia ordenada dentro del tratamiento disciplinado del corpus.
\end{observation}

\section{Condición mínima de nivel}
\begin{definition}
Dos formulaciones pertenecen a niveles distintos cuando no son intercambiables sin pérdida discernible de:
\begin{enumerate}
  \item dependencia estructural,
  \item capacidad operativa,
  \item función dentro del canon.
\end{enumerate}
\end{definition}

\begin{observation}
Esta cláusula evita confundir diferencia de longitud verbal con diferencia real de nivel.
\end{observation}

\section{Jerarquía y umbral}
\begin{definition}
Una formulación solo puede ocupar un nivel operativo explícito si ha atravesado un umbral suficiente para sostener tratamiento disciplinado en ese nivel.
\end{definition}

\begin{observation}
La jerarquía mínima depende del umbral operativo: sin él, el orden entre niveles corre el riesgo de ser nominal y no estructural.
\end{observation}

\section{Jerarquía y transición}
\begin{definition}
Toda transición de nivel disciplinada produce o explicita una relación jerárquica entre la formulación heredada y la formulación fortalecida.
\end{definition}

\begin{observation}
La transición no crea jerarquía de manera arbitraria; la vuelve legible cuando la dependencia y la ganancia operativa pueden declararse sin ambigüedad.
\end{observation}

\section{Condición mínima de fallo}
\begin{definition}
Se considera fallida una jerarquía declarada cuando los niveles propuestos resultan intercambiables, carecen de dependencia visible o no exhiben diferencia operativa discernible.
\end{definition}

\begin{observation}
Esta condición protege el canon intermedio frente a escalas infladas o puramente nominales.
\end{observation}

\section{Cláusula de alcance}
\begin{observation}
La jerarquía mínima fija una primera gramática de orden entre formulaciones dentro del Tomo II. No constituye todavía una taxonomía exhaustiva de todos los niveles, regímenes o dimensiones del sistema.
\end{observation}

